\reading{MR-1}{Estimating Market Risk Measures: An Introduction and Overview}{strike}{6}
% ==========================================================

\begin{mrkeybox}[title={Learning objectives in this reading}]
\small
\begin{enumerate}[label=\textbf{\alph*.},leftmargin=1.7em]
\item Estimate VaR using a historical simulation approach.
\item Estimate VaR using a parametric approach for both normal and lognormal return distributions.
\item Estimate the expected shortfall given profit and loss (P\&L) or return data.
\item Estimate risk measures by estimating quantiles.
\item Evaluate estimators of risk measures by estimating their standard errors.
\item Interpret quantile-quantile (QQ) plots to identify the characteristics of a distribution.
\end{enumerate}
\end{mrkeybox}

\section{Basics}

\subsection{Arithmetic returns}
This approach is not appropriate for long investment horizons and does not consider
reinvestment.

\begin{mrfmlbox}
\[ r_t \;=\; \frac{P_t + D_t}{P_{t-1}} \;-\; 1 \]
\mrvk{$P_t$ = asset price at the end of period $t$;\quad $D_t$ = interim payment (dividend or
coupon) received during period $t$;\quad $P_{t-1}$ = asset price at the start of the period.}{}
\end{mrfmlbox}

\subsection{Geometric returns}
Here, interim payments are continuously reinvested. Note that this approach ensures that
asset price can never be negative.

\begin{mrfmlbox}
\[ R_t \;=\; \ln\!\left( \frac{P_t + D_t}{P_{t-1}} \right) \]
\mrvk{$R_t$ = geometric (log) return over period $t$; the remaining symbols are as above.}{}
\end{mrfmlbox}

\subsection{Data}
Our data can come in various forms. Perhaps the simplest is in terms of profit/loss (or P/L).

% ---------------------------------------------------------
\lo{MR-1 a}{Estimate VaR using a historical simulation approach.}

Estimating VaR with a historical simulation approach is by far the simplest and most
straightforward VaR method.

The idea here is to order the historical losses observed in a portfolio, ordered by the actual
amount of the loss. To make this calculation, you simply \term{order return observations from
largest to smallest}. The observation that follows the threshold loss level denotes the VaR limit.

So, if we have 1,000 loss observations and we are interested in the 97\% confidence level,
that implies that 30 observations will be in the tail of the distribution that will define the
VaR of this portfolio.

\begin{mrfmlbox}
\[ \text{VaR observation number} \;=\; (\alpha \times n) + 1 \]
\mrvk{$\alpha$ = tail probability, i.e.\ $1 - \text{confidence level}$;\quad $n$ = number of
observations in the data set.}{}
\end{mrfmlbox}

So, in the example the \term{31st} loss observation would be the amount of the expected VaR of
the portfolio.

\begin{center}
\begin{tikzpicture}
\begin{axis}[
  width=11.4cm, height=6.2cm,
  ybar interval, bar width=1,
  xmin=-4, xmax=4, ymin=0, ymax=92,
  xlabel={\scriptsize Profit / loss}, ylabel={\scriptsize Frequency},
  xlabel style={font=\scriptsize}, ylabel style={font=\scriptsize},
  tick label style={font=\scriptsize},
  axis background/.style={fill=white},
  xtick={-4,-3,-2,-1,0,1,2,3,4},
  ytick={0,10,20,30,40,50,60,70,80,90},
  grid=none]
\addplot[fill=PaleNavy, draw=FRMNavy!60, line width=0.25pt] coordinates {
(-4.000,0) (-3.850,0) (-3.700,0) (-3.550,1) (-3.400,1) (-3.250,0) (-3.100,0)
(-2.950,1) (-2.800,2) (-2.650,5) (-2.500,6) (-2.350,8) (-2.200,10) (-2.050,11)
(-1.900,20) (-1.750,14) (-1.600,27) (-1.450,28) (-1.300,31) (-1.150,44)
(-1.000,50) (-0.850,53) (-0.700,57) (-0.550,73) (-0.400,85) (-0.250,82)
(-0.100,65) (0.050,71) (0.200,70) (0.350,73) (0.500,75) (0.650,61) (0.800,55)
(0.950,35) (1.100,43) (1.250,38) (1.400,29) (1.550,26) (1.700,20) (1.850,9)
(2.000,3) (2.150,8) (2.300,7) (2.450,0) (2.600,2) (2.750,0) (2.900,0)
(3.050,0) (3.200,1) (3.350,0) (3.500,0) (3.650,0) (3.800,0) (3.950,0)};
\draw[FRMRed, -{Stealth[length=5pt]}, line width=0.9pt]
  (axis cs:-1.704,68) -- (axis cs:-1.704,8);
\node[anchor=west, font=\scriptsize\sffamily, text=FRMRed]
  at (axis cs:-3.90,72) {95\% VaR $=1.704$};
\end{axis}
\end{tikzpicture}
\end{center}
\figcap{The full profit/loss distribution built from the historical observations. The marker sits
at the loss beyond which only the worst 5\% of observations fall, so everything to its left is
the tail that defines the 95\% VaR.}

\begin{mrtrapbox}[title={Problems with the HS approach}]
From a practical perspective, the historical simulation approach is sensible \textit{only} if
you expect future performance to follow the same return generating process as in the past.
Furthermore, this approach is unable to adjust for changing economic conditions or abrupt
shifts in parameter values.
\end{mrtrapbox}

% ---------------------------------------------------------
\lo{MR-1 b}{Estimate VaR using a parametric approach for both normal and lognormal return distributions.}

\term{Parametric approach:} in this method, we use some parameters (like mean and variance)
--- for example the delta-normal approach --- and also explicitly assume a distribution for the
underlying observations.

\subsection{1) Normal VaR}
\begin{itemize}
\item We assume that arithmetic returns are normally distributed with mean $\mu$ and standard
      deviation $\sigma$.
\item For returns that follow a normal distribution, we use normal VaR.
\item The VaR cutoff will be in the left tail of the returns distribution. Hence, the calculated
      value at risk is negative, but is typically reported as a positive value since the negative
      amount is implied.
\end{itemize}

\begin{mrfmlbox}
\[ \mathrm{VaR}(\alpha\%) \;=\; -\mu_{P/L} \;+\; \sigma_{P/L}\, z_\alpha
   \qquad\qquad
   \mathrm{VaR}(\alpha\%) \;=\; \bigl(-\mu_r + \sigma_r \times z_\alpha\bigr) \times P_{t-1} \]
\mrvk{$\mu_{P/L}$, $\sigma_{P/L}$ = mean and standard deviation of the profit/loss distribution
(used when the data is given in currency units);\quad $\mu_r$, $\sigma_r$ = mean and standard
deviation of arithmetic returns (used when the data is given in \%, so the result must be scaled
by the portfolio value);\quad $z_\alpha$ = one-tailed standard normal deviate for the chosen
confidence level;\quad $P_{t-1}$ = portfolio value at the start of the period.}{}
\end{mrfmlbox}

\begin{mrexbox}[title={Example 1 --- P/L data}]
Assume that the profit/loss distribution for XYZ is normally distributed with an annual mean of
\$15 million and a standard deviation of \$10 million. Calculate the VaR at the 95\% and 99\%
confidence levels using a parametric approach.
\tcblower
\small
\begin{align*}
\mathrm{VaR}(95\%) &= -15 + 10 \times 1.645 = \$1.45\text{ million}\\
\mathrm{VaR}(99\%) &= -15 + 10 \times 2.326 = \$8.26\text{ million}
\end{align*}
\end{mrexbox}

\begin{mrexbox}[title={Example 2 --- arithmetic return data}]
A portfolio has a beginning period value of \$100. The arithmetic returns follow a normal
distribution with a mean of 10\% and a standard deviation of 20\%. Calculate VaR at both the
95\% and 99\% confidence levels.
\tcblower
\small
\begin{align*}
\mathrm{VaR}(95\%) &= \bigl(-0.10 + 0.20 \times 1.645\bigr) \times 100 = \$22.90\\
\mathrm{VaR}(99\%) &= \bigl(-0.10 + 0.20 \times 2.326\bigr) \times 100 = \$36.52
\end{align*}
\end{mrexbox}

\subsection{2) Lognormal VaR}
\begin{itemize}
\item Assume that geometric returns are normally distributed with mean $\mu$ and standard
      deviation $\sigma$. This assumption implies that the natural logarithm of $p$ is normally
      distributed, or that $p$ itself is lognormally distributed. Normally distributed geometric
      returns imply that the VaR is lognormally distributed.
\item The lognormal distribution is right-skewed with positive outliers and bounded below by
      zero. As a result, the lognormal distribution is commonly used to counter the possibility
      of negative asset prices ($P_t$).
\item Used when geometric returns are provided in the question.
\end{itemize}

\begin{mrfmlbox}
\[ \mathrm{VaR}(\alpha\%) \;=\; P_{t-1} \times \left( 1 - e^{\,\mu_R - \sigma_R \times z_\alpha} \right) \]
\mrvk{$\mu_R$, $\sigma_R$ = mean and standard deviation of \textit{geometric} returns;\quad
$z_\alpha$ = one-tailed standard normal deviate;\quad $P_{t-1}$ = portfolio value at the start
of the period.}{}
\end{mrfmlbox}

\begin{mrexbox}[title={Example 3 --- lognormal VaR}]
A diversified portfolio exhibits a normally distributed geometric return with mean and standard
deviation of 10\% and 20\%, respectively. Calculate the 5\% and 1\% lognormal VaR assuming the
beginning period portfolio value is \$100.
\tcblower
\small
\begin{align*}
\mathrm{VaR}(5\%) &= 100 \times \left(1 - e^{\,0.10 - 0.20 \times 1.645}\right)
   = 100 \times \left(1 - e^{-0.229}\right) = \$20.47\\
\mathrm{VaR}(1\%) &= 100 \times \left(1 - e^{\,0.10 - 0.20 \times 2.326}\right)
   = 100 \times \left(1 - e^{-0.3652}\right) = \$30.59
\end{align*}
\end{mrexbox}

\subsection{Interpretations of VaR}

\begin{mrdefbox}[title={1. Standard interpretation --- VaR is a loss}]
Assume a risk manager calculates Normal 99\% VaR as \$5,000.
\begin{itemize}
\item \term{Confidence level:} I am 99 percent confident that the maximum 1-week loss will not
      exceed \$5,000.
\item \term{Minimum loss level:} there is a 1 percent chance that the loss will exceed \$5,000
      or more.
\end{itemize}
\end{mrdefbox}

\begin{mrdefbox}[title={2. Alternate interpretation --- VaR is negative, so it is a profit}]
Assume a risk manager calculates Normal 99\% VaR as \textsc{negative} \$5,000. When we have
profits, both the minimum loss and the confidence interval statement are reversed.
\begin{itemize}
\item \term{Confidence level:} I am 99 percent confident that the minimum profit will be at
      least \$5,000.
\item \term{Minimum loss level:} there is a 1 percent chance that the profit will be less than
      \$5,000.
\end{itemize}
\end{mrdefbox}

\begin{mrexbox}[title={Example 4 --- interpreting a negative VaR}]
Assume the profit/loss distribution for XYZ is normally distributed with an annual mean of
\$20 million and a standard deviation of \$10 million. The 5\% VaR is calculated and interpreted
as which of the following statements?
\begin{enumerate}[label=\Alph*.,leftmargin=1.7em]
\item 5\% probability of losses of at least \$3.50 million.
\item 5\% probability of earnings of at least \$3.50 million.
\item 95\% probability of losses of at least \$3.50 million.
\item 95\% probability of earnings of at least \$3.50 million.
\end{enumerate}
\tcblower
\small
$\mathrm{VaR}(5\%) = -20 + 10 \times 1.65 = -\$3.50$ million. The VaR is negative, so it is a
profit, and the alternate interpretation applies. \term{Answer: B.}
\end{mrexbox}

% ---------------------------------------------------------
\lo{MR-1 c}{Estimate the expected shortfall given profit and loss (P\&L) or return data.}

\begin{mrtrapbox}[title={Limitations of VaR}]
\begin{itemize}
\item Confidence level dependent.
\item Does not tell the magnitude of the actual loss.
\end{itemize}
\end{mrtrapbox}

\begin{mrdefbox}[title={Expected shortfall}]
The expected shortfall (ES) provides an estimate of the tail loss by averaging the VaRs for
increasing confidence levels in the tail.
\end{mrdefbox}

\subsection{Two methods to calculate ES}

\textbf{1) Slicing the tail.} Specifically, the tail mass is divided into $n$ equal slices and
the corresponding $n-1$ VaRs are computed.

\medskip
\begin{center}
\small
\begin{tabular}{>{\raggedright\arraybackslash}p{4.2cm} r r}
\arrayrulecolor{FRMRule}\toprule
\rowcolor{FRMNavy} \hdr{Confidence level} & \hdr{VaR} & \hdr{Difference}\\
96\% & 1.7507 & \\
\rowcolor{FRMSoft} 97\% & 1.8808 & 0.1301\\
98\% & 2.0537 & 0.1729\\
\rowcolor{FRMSoft} 99\% & 2.3263 & 0.2726\\
\midrule
\textit{Average} & 2.003 & \\
\textit{Theoretical true value} & 2.063 & \\
\bottomrule
\end{tabular}
\end{center}

Observe that the VaR increases (from the Difference column) in order to maintain the same
interval mass (of 1\%) because the tails become thinner and thinner.

The average of the four computed VaRs is \term{2.003} and represents the probability-weighted
expected tail loss (a.k.a.\ expected shortfall). Note that as $n$ increases, the expected
shortfall will increase and approach the theoretical true loss [2.063 in this case; the average
of a high number of VaRs (e.g., greater than 10,000)].

\medskip
\begin{center}
\small
\begin{tabular}{>{\raggedright\arraybackslash}p{5.0cm} r}
\arrayrulecolor{FRMRule}\toprule
\rowcolor{FRMNavy} \hdr{Number of tail slices ($n$)} & \hdr{ES}\\
10 & 2.0250\\
\rowcolor{FRMSoft} 25 & 2.0433\\
50 & 2.0513\\
\rowcolor{FRMSoft} 100 & 2.0562\\
250 & 2.0597\\
\rowcolor{FRMSoft} 500 & 2.0610\\
1,000 & 2.0618\\
\rowcolor{FRMSoft} 2,500 & 2.0623\\
5,000 & 2.0625\\
\rowcolor{FRMSoft} 10,000 & 2.0626\\
\midrule
\textit{True value} & 2.0630\\
\bottomrule
\end{tabular}
\end{center}

\textbf{2) Direct formula.} The second method is using this formula to calculate ES.

\begin{mrfmlbox}
\[ \mathrm{ES} \;=\; \mu \;+\; \sigma\,\frac{e^{-(z^2/2)}}{(1-x)\sqrt{2\pi}} \]
\mrvk{$\mu$, $\sigma$ = mean and standard deviation of the distribution;\quad
$z$ = the standard normal quantile at the chosen confidence level;\quad
$x$ = the confidence level itself, so $(1-x)$ is the tail probability.}{%
\par\smallskip\textbf{\sffamily Note on notation:} $z$ and $x$ are two different things in the
same formula --- $z$ sits in the numerator as the quantile, $x$ sits in the denominator as the
confidence level.}
\end{mrfmlbox}

% ---------------------------------------------------------
\lo{MR-1 d}{Estimate risk measures by estimating quantiles.}

Under expected shortfall estimation, the tail region is divided into equal probability slices
and then multiplied by the corresponding quantiles.

\begin{center}
\begin{tikzpicture}
\begin{axis}[
  width=11.0cm, height=5.0cm,
  ybar, bar width=3.2pt,
  xmin=-4.2, xmax=4.2, ymin=0, ymax=0.45,
  axis y line=none, axis x line=bottom,
  axis line style={draw=black!55},
  tick label style={font=\scriptsize}, xtick={-4,-3,-2,-1,0,1,2,3,4},
  x axis line style={-},
  title={\small\sffamily\bfseries\color{FRMGold}Expected Shortfall},
  clip=false]
\addplot[fill=FRMRed!45, draw=FRMRed!70, line width=0.2pt] coordinates {
(-3.4,0.0012)(-3.2,0.0024)(-3.0,0.0044)(-2.8,0.0079)(-2.6,0.0136)(-2.4,0.0224)
(-2.2,0.0355)(-2.0,0.0540)(-1.8,0.0790)};
\addplot[fill=PaleNavy, draw=FRMNavy!60, line width=0.2pt] coordinates {
(-1.6,0.1109)(-1.4,0.1497)(-1.2,0.1942)(-1.0,0.2420)(-0.8,0.2897)(-0.6,0.3332)
(-0.4,0.3683)(-0.2,0.3910)(0.0,0.3989)(0.2,0.3910)(0.4,0.3683)(0.6,0.3332)
(0.8,0.2897)(1.0,0.2420)(1.2,0.1942)(1.4,0.1497)(1.6,0.1109)(1.8,0.0790)
(2.0,0.0540)(2.2,0.0355)(2.4,0.0224)(2.6,0.0136)(2.8,0.0079)(3.0,0.0044)
(3.2,0.0024)(3.4,0.0012)};
\addplot[domain=-4:4, samples=200, smooth, black!80, line width=0.9pt, forget plot]
  {0.39894*exp(-x^2/2)};
\node[anchor=south east, font=\scriptsize\sffamily, align=right, text=black!70]
  at (axis cs:-2.25,0.20) {Equal weight to\\all these losses};
\draw[black!70, -{Stealth[length=4pt]}, line width=0.5pt]
  (axis cs:-2.35,0.195) -- (axis cs:-2.75,0.045);
\draw[black!70, -{Stealth[length=4pt]}, line width=0.5pt]
  (axis cs:-2.20,0.195) -- (axis cs:-1.90,0.10);
\end{axis}
\end{tikzpicture}
\end{center}
\figcap{The shaded tail is cut into equal-probability slices. A VaR is taken at each slice
boundary and the slices are then averaged with equal weight, which is what makes the result an
expected shortfall rather than a single quantile.}

Risk measure estimation by quantiles means to divide the entire distribution into different
quantiles ($n$).

\begin{mrdefbox}[title={Coherent risk measure}]
\begin{itemize}
\item A more general risk measure than either VaR or ES is known as a \term{coherent risk
      measure}. A coherent risk measure is a weighted average of the quantiles of the loss
      distribution where the weights are user-specific based on individual risk aversion.
\item A coherent risk measure will assign \textit{each} quantile (not just tail quantiles) a
      weight. The average of the weighted VaRs is the estimated loss.
\end{itemize}
\end{mrdefbox}

% ---------------------------------------------------------
\lo{MR-1 e}{Evaluate estimators of risk measures by estimating their standard errors.}

\term{Standard error:} used to assess the precision of the risk measure. In Value at Risk (VaR)
it refers to the uncertainty or the variability of the VaR estimate. The simplest method is to
create a confidence interval around the quantile in question.

Estimators that are less precise (i.e., have large standard errors and wide confidence
intervals) will have limited practical value. Therefore, it is best practice to also compute
the standard error.

\begin{mrfmlbox}[breakable=false]
\textbf{a. Confidence interval}
\[ \bigl[\, q + se(q) \times z_\alpha \,\bigr] \;>\; \mathrm{VaR} \;>\;
   \bigl[\, q - se(q) \times z_\alpha \,\bigr] \]
\mrvk{$q$ = the estimated quantile;\quad $se(q)$ = standard error of that quantile;\quad
$z_\alpha$ = the VaR $z$-value, one tail.}{}

\medskip
\textbf{b.} $f(q)$ is the area of a bin width, $n$ is sample size, $p$ is tail probability.

\medskip
\textbf{c. Standard error}
\[ se(q) \;=\; \frac{\sqrt{p\,(1-p)\,/\,n}}{f(q)} \]
\end{mrfmlbox}

\begin{mrkeybox}[title={How to reduce standard errors}]
\begin{enumerate}
\item Increase sample size ($n$) $\Rightarrow$ SE down.
\item Increase bin size $f(q)$ $\Rightarrow$ SE down.
\item Reduce tail probability ($p$) $\Rightarrow$ SE down.
\end{enumerate}
\end{mrkeybox}

% ---------------------------------------------------------
\lo{MR-1 f}{Interpret quantile-quantile (QQ) plots to identify the characteristics of a distribution.}

\term{Quantile-quantile plots.} The quantile-quantile (QQ) plot is a straightforward way to
visually examine if empirical data fits the reference or hypothesized theoretical distribution.
It is a plot of the quantiles of the empirical distribution against those of some specified
distribution. The shape of the QQ plot tells us a lot about how the empirical distribution
compares to the specified one.

\begin{itemize}
\item In particular, if the QQ plot is \term{linear}, then the specified distribution fits the
      data, and we have identified the distribution to which our data belong.
\end{itemize}

\begin{center}
\begin{tikzpicture}
% ===== ROW 1: linear QQ -> the reference density fits =====
\begin{axis}[name=qq1, width=5.7cm, height=4.1cm,
  xmin=-4,xmax=4,ymin=-4,ymax=4,
  xlabel={\scriptsize Normal quantiles}, ylabel={\scriptsize Empirical quantiles},
  xlabel style={font=\tiny}, ylabel style={font=\tiny},
  tick label style={font=\tiny}, xtick={-2,0,2}, ytick={-2,0,2}, grid=none,
  title={\scriptsize\sffamily\bfseries\color{FRMNavy}QQ plot is \textit{linear}},
  axis line style={draw=black!55}]
\addplot[only marks, mark=+, mark size=1.2pt, color=black!55, domain=-3.2:3.2, samples=70]
  {x + 0.05*sin(deg(3*x))};
\addplot[dash dot, black!60, line width=0.6pt, domain=-3.4:3.4, forget plot] {x};
\end{axis}
\draw[FRMGold, -{Stealth[length=5.5pt]}, line width=1.2pt]
  ($(qq1.east)+(0.55cm,0)$) -- ($(qq1.east)+(1.75cm,0)$);
\node[anchor=south, font=\tiny\sffamily, text=FRMGold, align=center]
  at ($(qq1.east)+(1.15cm,0.12cm)$) {implies};
\begin{axis}[name=d1, at={($(qq1.east)+(2.0cm,0)$)}, anchor=west,
  width=5.7cm, height=4.1cm,
  xmin=-4.2, xmax=4.2, ymin=0, ymax=0.46,
  axis y line=none, axis x line=bottom, axis line style={draw=black!55},
  xlabel={\scriptsize Return}, xlabel style={font=\tiny},
  xtick={-2,0,2}, tick label style={font=\tiny},
  title={\scriptsize\sffamily\bfseries\color{FRMNavy}the data \textit{is} normal},
  clip=false]
\addplot[FRMNavy, line width=1.2pt, smooth, domain=-4.2:4.2, samples=140]
  {0.39894*exp(-(x^2)/2)};
\addplot[black!55, line width=0.9pt, dashed, smooth, domain=-4.2:4.2, samples=140]
  {0.39894*exp(-(x^2)/2)};
\node[anchor=north, font=\tiny\sffamily, text=black!60] at (axis cs:2.55,0.115)
  {curves coincide};
\end{axis}
% ===== ROW 2: S-shaped QQ -> fatter tails =====
\begin{axis}[name=qq2, at={($(qq1.south)+(0,-1.85cm)$)}, anchor=north,
  width=5.7cm, height=4.1cm,
  xmin=-4,xmax=4,ymin=-8,ymax=8,
  xlabel={\scriptsize Normal quantiles}, ylabel={\scriptsize Empirical quantiles},
  xlabel style={font=\tiny}, ylabel style={font=\tiny},
  tick label style={font=\tiny}, xtick={-2,0,2}, ytick={-5,0,5}, grid=none,
  title={\scriptsize\sffamily\bfseries\color{FRMRed}QQ plot bends at \textit{both} ends},
  axis line style={draw=black!55}]
\addplot[only marks, mark=+, mark size=1.2pt, color=black!55, domain=-3.1:3.1, samples=70]
  {x + 0.14*x^3};
\addplot[dash dot, black!60, line width=0.6pt, domain=-3.2:3.2, forget plot] {1.8*x};
\end{axis}
\draw[FRMGold, -{Stealth[length=5.5pt]}, line width=1.2pt]
  ($(qq2.east)+(0.55cm,0)$) -- ($(qq2.east)+(1.75cm,0)$);
\node[anchor=south, font=\tiny\sffamily, text=FRMGold, align=center]
  at ($(qq2.east)+(1.15cm,0.12cm)$) {implies};
\begin{axis}[name=d2, at={($(qq2.east)+(2.0cm,0)$)}, anchor=west,
  width=5.7cm, height=4.1cm,
  xmin=-4.2, xmax=4.2, ymin=0, ymax=0.46,
  axis y line=none, axis x line=bottom, axis line style={draw=black!55},
  xlabel={\scriptsize Return}, xlabel style={font=\tiny},
  xtick={-2,0,2}, tick label style={font=\tiny},
  title={\scriptsize\sffamily\bfseries\color{FRMRed}the data has \textit{fatter tails}},
  clip=false]
\addplot[FRMRed, line width=1.2pt, smooth, domain=-4.2:4.2, samples=180]
  {0.425*exp(-(x^2)/1.15) + 0.048*exp(-(x^2)/16)};
\addplot[black!55, line width=0.9pt, dashed, smooth, domain=-4.2:4.2, samples=140]
  {0.39894*exp(-(x^2)/2)};
\node[anchor=south, font=\tiny\sffamily, text=FRMRed, fill=white, inner sep=1pt]
  at (axis cs:-3.1,0.055) {heavier};
\node[anchor=south, font=\tiny\sffamily, text=FRMRed, fill=white, inner sep=1pt]
  at (axis cs:3.1,0.055) {heavier};
\end{axis}
\end{tikzpicture}
\end{center}
\figcap{The QQ plot on the left and the density on the right carry the same information. Top row:
the points sit on the reference line, so the empirical distribution matches the hypothesised one
and the two densities coincide. Bottom row: the points climb \textit{above} the line at the right
end and fall \textit{below} it at the left end. That bend means the empirical quantiles are more
extreme than the reference at both ends, which is exactly what a fatter-tailed density looks
like. Reading the bend is how you diagnose the tail without ever plotting the density.}

% ==========================================================
